% =====================================================================
\section{Le front en détail}

\subsection{Rien de construit, tout d'écrit}

Le projet n'utilise \textbf{aucun \emph{framework} front} : pas de React,
pas de Bootstrap, pas d'étape de compilation JavaScript. Le HTML est produit
par le serveur, la mise en forme tient dans une feuille de style écrite à la
main, et deux bibliothèques seulement sont chargées depuis un CDN :

\noindent\begin{tabularx}{\textwidth}{|>{\columncolor{btkpale}\bfseries}L{3.4cm}|X|}
\hline
ApexCharts & \texttt{cdn.jsdelivr.net/npm/apexcharts} --- les trois graphiques du tableau de bord\\
\hline
Lucide & \texttt{unpkg.com/lucide@latest} --- les icônes (\texttt{<i data-lucide="\dots">})\\
\hline
\end{tabularx}

\begin{retenir}
\textbf{Pourquoi ce choix ?} Une application interne, rendue côté serveur,
n'a pas besoin d'un client riche : les pages sont des listes et des
formulaires. Ne pas ajouter de \emph{framework}, c'est aussi ne pas ajouter
d'étape de construction, de gestionnaire de paquets ni de dépendances à
maintenir. La contrepartie est qu'il n'y a pas de composants réutilisables :
le facteur commun est obtenu par les inclusions JSP (\texttt{sidebar},
\texttt{top-bar}) et par les classes CSS.
\end{retenir}

\subsection{La feuille de style}

\texttt{css/style.css}, 794 lignes, organisée autour de \textbf{variables CSS}
déclarées sur \texttt{:root} --- couleurs de la charte BTK, fonds, texte,
états. Changer la charte se fait en un endroit. Les classes principales sont
\texttt{card}, \texttt{stat-card-modern}, \texttt{nav-item}, \texttt{btn},
\texttt{form-control}, \texttt{table}.

\subsection{JSTL et le langage d'expression}

Les vues n'embarquent \textbf{aucun code Java} : pas de scriptlet
\texttt{<\% \dots \%>}. Tout passe par JSTL et le langage d'expression, ce qui
garde la vue lisible et sans logique métier.

\noindent\begin{tabularx}{\textwidth}{|>{\ttfamily}L{4.6cm}|X|}
\hline
\textmd{\textbf{Balise}} & \textbf{Usage dans le projet}\\
\hline
<c:if test="\$\{\dots\}"> & afficher un bloc selon le rôle ou l'état\\
\hline
<c:forEach var items> & parcourir une liste posée par le servlet\\
\hline
<c:choose> <c:when> & alternative à plusieurs branches\\
\hline
<c:set var value> & variable locale de vue\\
\hline
\$\{sessionScope.role\} & lire un attribut de session\\
\hline
\$\{pageContext.request.\newline contextPath\} & préfixer les liens par \texttt{/gestion-agences}\\
\hline
\end{tabularx}

\subsection{Les trois graphiques du tableau de bord}

Le servlet ne fabrique pas de JSON : il pose des \textbf{chaînes déjà
formatées}, que la JSP insère dans le script.

\begin{lstlisting}[style=java]
// HomeServlet
request.setAttribute("chartPointageLabels", "'Presents','Retards','Absents'");
request.setAttribute("chartPointageData", presentsStricts + "," + retards + "," + absents);
\end{lstlisting}

\begin{lstlisting}[style=jsp]
new ApexCharts(document.querySelector("#pointageChart"), {
    chart: { type: 'donut', height: 320, background: 'transparent' },
    labels: [${chartPointageLabels}],
    series: [${chartPointageData}]
}).render();
\end{lstlisting}

C'est simple et sans dépendance de sérialisation. La limite est qu'un libellé
contenant une apostrophe casserait le script : pour des données saisies par
l'utilisateur, il faudrait passer par du JSON échappé.

% =====================================================================
\section{Le back en détail}

\subsection{Le motif commun à tous les servlets}

Les dix-huit servlets suivent la même structure, ce qui rend le code
prévisible :

\begin{enumerate}[leftmargin=*,itemsep=1pt]
  \item récupérer la session avec \texttt{getSession(false)} ;
  \item contrôler l'authentification, puis le rôle ;
  \item lire les paramètres de la requête ;
  \item appeler un ou plusieurs services ;
  \item \texttt{setAttribute} puis \texttt{forward} pour afficher, ou
        \texttt{sendRedirect} après une écriture.
\end{enumerate}

\subsection{Les EJB}

Quinze services sont \texttt{@Stateless}, un seul est \texttt{@Singleton}.

\noindent\begin{tabularx}{\textwidth}{|>{\columncolor{btkpale}\bfseries}L{3.2cm}|X|}
\hline
\texttt{@Stateless} & Le conteneur gère un réservoir d'instances et en prête une le temps de l'appel. Aucun état conservé entre deux requêtes, donc pas de problème de concurrence.\\
\hline
\texttt{@Singleton} & Une seule instance pour toute l'application. Utilisé pour \texttt{SecurityService}, avec \texttt{@Startup} pour qu'elle existe dès le démarrage.\\
\hline
\end{tabularx}

\vspace{0.5em}
L'injection est déclarative --- aucun \texttt{new} sur un service :

\begin{lstlisting}[style=java]
@EJB
private DemandeModificationService demandeModificationService;
\end{lstlisting}

\subsection{JPA et les transactions}

Chaque service reçoit son \texttt{EntityManager} par injection :

\begin{lstlisting}[style=java]
@PersistenceContext(unitName = "default")
private EntityManager em;
\end{lstlisting}

Les requêtes sont écrites en \textbf{JPQL} --- sur les entités, pas sur les
tables --- et \textbf{toujours paramétrées} :

\begin{lstlisting}[style=java]
return em.createQuery(
        "SELECT d FROM DemandeModification d WHERE d.utilisateur = :u "
      + "ORDER BY d.dateDemande DESC", DemandeModification.class)
    .setParameter("u", user)
    .getResultList();
\end{lstlisting}

\begin{retenir}
\textbf{Question quasi certaine : et l'injection SQL ?} Aucune requête n'est
construite par concaténation de chaînes ; les valeurs passent par
\texttt{setParameter}, donc par une variable liée côté pilote. La valeur n'est
jamais interprétée comme du SQL.
\end{retenir}

\textbf{Les transactions ne sont pas écrites.} L'unité de persistance est
déclarée \texttt{transaction-type="JTA"} et les services sont des EJB : le
conteneur ouvre une transaction à l'entrée de chaque méthode et la valide à la
sortie, ou l'annule si une exception s'échappe. C'est la gestion \emph{par le
conteneur} (CMT). D'où, dans la trace d'erreur de connexion, la présence de
\texttt{CMTTxInterceptor} : c'est l'intercepteur qui pose la transaction.

\subsection{L'unité de persistance}

\begin{lstlisting}[style=xml]
<persistence-unit name="default" transaction-type="JTA">
    <jta-data-source>java:/OracleDS</jta-data-source>
    ... 16 classes d'entites ...
    <properties>
        <property name="hibernate.dialect" value="org.hibernate.dialect.OracleDialect"/>
        <property name="hibernate.hbm2ddl.auto" value="none"/>
    </properties>
</persistence-unit>
\end{lstlisting}

Deux points méritent d'être signalés au jury :

\begin{itemize}[leftmargin=*,itemsep=2pt]
  \item \texttt{hbm2ddl.auto = none} : Hibernate ne crée ni ne modifie
        \textbf{jamais} le schéma. Les tables sont créées par les scripts de
        \texttt{sql/}, et la base fait autorité.
  \item aucun identifiant de connexion n'est écrit ici. Ils appartiennent à la
        datasource \texttt{java:/OracleDS}, configurée dans WildFly. Le même
        WAR se déploie donc sur plusieurs environnements sans être recompilé.
\end{itemize}
